% Part: second-order-logic
% Chapter: syntax-and-semantics
% Section: satisfaction

\documentclass[../../../include/open-logic-section]{subfiles}

\begin{document}

\olfileid{sol}{syn}{sat}

\olsection{સંતોષ}

\begin{explain}
દ્વિતીય-ક્રમ !!{formula}s માટે સંતોષ-સંબંધ~$\Sat{M}{!A}[s]$
વ્યાખ્યાયિત કરવા માટે આપણે દ્વિતીય-ક્રમ !!{variable}sને પણ આવરી લે
એ રીતે અગાઉની વ્યાખ્યાઓ વિસ્તૃત કરવી પડે. !!a{structure}નો ખ્યાલ
દ્વિતીય-ક્રમ અને પ્રથમ-ક્રમ તર્કશાસ્ત્રમાં એકસરખો છે. ફરક માત્ર
ચલ-નિયુક્તિઓ~$s$માં છે: હવે તેમણે પ્રથમ-ક્રમ !!{variable}s ઉપરાંત
દ્વિતીય-ક્રમ !!{variable}sને પણ મૂલ્યો આપવા પડે.
\end{explain}

\begin{defn}[ચલ-નિયુક્તિ]
!!a{structure}~$\Struct{M}$ માટેની \emph{ચલ-નિયુક્તિ}~$s$ એવું
વિધેય છે જે દરેક
\begin{enumerate}
\item વસ્તુ-!!{variable}~$\Obj{v_i}$ને~$\Domain M$નો ઘટક આપે છે,
  એટલે $s(\Obj{v_i}) \in \Domain{M}$;
\item $n$-સ્થાની સંબંધ-ચલ~$\Obj{V_i^n}$ને~$\Domain{M}$ પરનો
  $n$-સ્થાની સંબંધ આપે છે, એટલે
  $s(\Obj{V_i^n}) \subseteq \Domain{M}^n$;
\item $n$-સ્થાની વિધેય-ચલ~$\Obj{u_i^n}$ને $\Domain{M}$માંથી
  $\Domain{M}$માં જતું $n$-સ્થાની વિધેય આપે છે, એટલે
  $s(\Obj{u_i^n})\colon \Domain{M}^n \to \Domain{M}$.
\end{enumerate}
\end{defn}

\begin{explain}
!!^a{structure} દરેક !!{constant} અને !!{function}ને !!a{value}
આપે છે; દ્વિતીય-ક્રમ ચલ-નિયુક્તિ દરેક વસ્તુ-ચલને વસ્તુ, દરેક
વિધેય-ચલને વિધેય અને દરેક સંબંધ-ચલને સંબંધ આપે છે. આ બંને સાથે
મળીને દરેક પદનું મૂલ્ય નક્કી કરી શકાય છે.
\sourcecorrection{OLSOL-002}{સ્થિર અંગ્રેજી મૂળની આ સમજણમાં
ચલ-નિયુક્તિ દ્વારા વસ્તુ-ચલો અને વિધેય-ચલોને મળતા મૂલ્યો જણાવ્યાં
છે, પરંતુ ઉપરની વ્યાખ્યામાં આવરી લીધેલા સંબંધ-ચલોને મળતા સંબંધોનો
ઉલ્લેખ રહી ગયો છે. અહીં તે ત્રીજો પ્રકાર પણ સ્પષ્ટ કર્યો છે.}
\end{explain}

\begin{defn}[પદનું \usetoken{S}{value}]
જો~$t$ ભાષા~$\Lang L$નું પદ હોય, $\Struct M$ એ~$\Lang L$ માટેની
!!a{structure} હોય અને~$s$ એ~$\Struct M$ માટેની !!a{variable}
નિયુક્તિ હોય, તો \emph{!!{value}}~$\Value{t}{M}[s]$ પ્રથમ-ક્રમ
પદોની જેમ જ, નીચેની વધારાની કલમ સાથે વ્યાખ્યાયિત થાય છે:
\begin{quote}
\indcase{t}{\Atom{u}{t_1, \ldots, t_n}}{
\[
\Value{\indfrm}{M}[s] = s(u)(\Value{t_1}{M}[s], \ldots,
\Value{t_n}{M}[s]).
\]}
\end{quote}
\end{defn}

\begin{defn}[$x$-રૂપાંતર]
જો~$s$ એ !!a{structure}~$\Struct M$ માટેની !!a{variable} નિયુક્તિ
હોય, તો~$\Struct M$ માટેની એવી કોઈપણ !!{variable} નિયુક્તિ~$s'$,
જે~$s$થી વધુમાં વધુ~$x$ને આપેલા મૂલ્યમાં જ જુદી હોય, તેને~$s$નું
\emph{$x$-રૂપાંતર} કહે છે. $s'$ એ~$s$નું $x$-રૂપાંતર હોય તો
$\varAssign{s'}{s}{x}$ લખીએ છીએ. (દ્વિતીય-ક્રમ ચલો~$X$ અથવા~$u$
માટે પણ આવું જ છે.)
\end{defn}

\begin{defn}
  જો~$s$ એ !!a{structure}~$\Struct M$ માટેની !!a{variable} નિયુક્તિ
  હોય અને~$m \in \Domain{M}$ હોય, તો~$\Subst{s}{m}{x}$ નીચે મુજબ
  વ્યાખ્યાયિત !!{variable} નિયુક્તિ છે:
  \[\Subst{s}{m}{y} = \begin{cases}
    m & \text{જો } y \ident x\\
    s(y) & \text{અન્યથા},
  \end{cases}\]
  જો~$X$ એ $n$-સ્થાની સંબંધ-!!{variable} હોય અને $M \subseteq
  \Domain{M}^n$ હોય, તો~$\Subst{s}{M}{X}$ નીચે મુજબ વ્યાખ્યાયિત
  !!{variable} નિયુક્તિ છે:
  \[\Subst{s}{M}{y} = \begin{cases}
    M & \text{જો } y \ident X\\
    s(y) & \text{અન્યથા}.
  \end{cases}\]
  જો~$u$ એ $n$-સ્થાની વિધેય-!!{variable} હોય અને $f\colon \Domain{M}^n
  \to \Domain{M}$ હોય, તો~$\Subst{s}{f}{u}$ નીચે મુજબ વ્યાખ્યાયિત
  !!{variable} નિયુક્તિ છે:
  \[\Subst{s}{f}{y} = \begin{cases}
    f & \text{જો } y \ident u\\
    s(y) & \text{અન્યથા}.
  \end{cases}\]
  દરેક કિસ્સામાં~$y$ કોઈપણ પ્રથમ-ક્રમ અથવા દ્વિતીય-ક્રમ !!{variable}
  હોઈ શકે છે.
\end{defn}

\begin{defn}[સંતોષ]
દ્વિતીય-ક્રમ !!{formula}s~$!A$ માટે સંતોષની વ્યાખ્યા
\olref[fol][syn][sat]{defn:satisfaction} જેવી જ છે, પરંતુ તેમાં
નીચેની કલમો ઉમેરાય છે:
\begin{enumerate}
\item \indcase{!A}{\Atom{X^n}{t_1, \dots, t_n}}{$\Sat{M}{\indfrm}[s]$
  ત્યારે અને માત્ર ત્યારે જ જ્યારે
  $\langle \Value{t_1}{M}[s], \dots, \Value{t_n}{M}[s] \rangle \in
  s(X^n)$.}
\tagitem{prvAll}{%
  \indcase{!A}{\lforall[X][!B]}{$\Sat{M}{\indfrm}[s]$ ત્યારે અને
  માત્ર ત્યારે જ જ્યારે દરેક $M \subseteq \Domain{M}^n$ માટે
  $\Sat{M}{!B}[\Subst{s}{M}{X}]$.}}{}

\tagitem{prvEx}{%
  \indcase{!A}{\lexists[X][!B]}{$\Sat{M}{\indfrm}[s]$ ત્યારે અને
  માત્ર ત્યારે જ જ્યારે ઓછામાં ઓછા એક $M \subseteq \Domain{M}^n$
  માટે $\Sat{M}{!B}[\Subst{s}{M}{X}]$.}}{}

\tagitem{prvAll}{%
  \indcase{!A}{\lforall[u][!B]}{$\Sat{M}{\indfrm}[s]$ ત્યારે અને
  માત્ર ત્યારે જ જ્યારે દરેક
  $f\colon \Domain{M}^n \to \Domain{M}$ માટે
  $\Sat{M}{!B}[\Subst{s}{f}{u}]$.}}{}

\tagitem{prvEx}{%
  \indcase{!A}{\lexists[u][!B]}{$\Sat{M}{\indfrm}[s]$ ત્યારે અને
  માત્ર ત્યારે જ જ્યારે ઓછામાં ઓછા એક
  $f\colon \Domain{M}^n \to \Domain{M}$ માટે
  $\Sat{M}{!B}[\Subst{s}{f}{u}]$.}}{}
\end{enumerate}
\end{defn}

\begin{ex}
  !!{formula}~$\lforall[z][(\Atom{X}{z} \liff \lnot
    \Atom{Y}{z})]$ વિચારો. તેમાં દ્વિતીય-ક્રમ પરિમાણકો નથી, પરંતુ
  દ્વિતીય-ક્રમ !!{variable}s~$X$ અને~$Y$ છે (અહીં બંને એકસ્થાની
  સમજવાના છે). તેને અનુરૂપ પ્રથમ-ક્રમ !!{sentence}
  $\lforall[z][(\Atom{P}{z} \liff \lnot \Atom{R}{z})]$ કહે છે કે
  જે કશું~$P$ના અર્થઘટનમાં આવે છે તે~$R$ના અર્થઘટનમાં આવતું નથી,
  અને ઊલટું પણ. કોઈ !!a{structure}માં !!a{predicate}~$P$નું અર્થઘટન
  $\Assign{P}{M}$ વડે મળે છે. પરંતુ~$X$ અને~$Y$ જેવા દ્વિતીય-ક્રમ
  !!{variable}sનું અર્થઘટન !!{structure} પોતે નહિ, પરંતુ
  !!a{variable} નિયુક્તિ આપે છે. આ દ્વિતીય-ક્રમ !!{formula},
  !!a{sentence} નથી (તેમાં મુક્ત !!{variable}s~$X$ અને~$Y$ છે), તેથી
  તેનો સંતોષ !!a{structure}~$\Struct{M}$ અને તેની સાથેની
  !!a{variable} નિયુક્તિ~$s$ને સાપેક્ષ જ નક્કી થાય છે.

  વ્યાપકક્ષેત્રના બધા ઘટકોને ધ્યાનમાં લઈએ તો, $s(X)$ના !!{element}s
  ચોક્કસ એ જ છે જે $s(Y)$ના !!{element}s નથી; એટલે દરેક ઘટક માટે આ
  દ્વિમુખી શરત સાચી હોય,
  ત્યારે $\Sat{M}{\lforall[z][(\Atom{X}{z} \liff \lnot
  \Atom{Y}{z})]}[s]$ સાચું છે; એટલે તે ત્યારે અને માત્ર ત્યારે જ
  જ્યારે $s(Y) = \Domain{M} \setminus s(X)$.
  \sourcecorrection{OLSOL-003}{સ્થિર અંગ્રેજી મૂળ અહીં બંને ગણોના
  ઘટકો એકબીજામાં નથી એવું કહે છે, જે માત્ર વિચ્છેદતા બતાવે છે;
  સૂત્રની દ્વિમુખી શરત મુજબ બંને ગણ પરસ્પર પૂરક હોવા જરૂરી છે.
  અહીં દરેક ઘટક માટે તે શરત સ્પષ્ટ કરી છે.}
  દાખલા તરીકે,
  $\Domain{M} = \{1, 2, 3\}$ લો. અહીં કોઈ !!{predicate}s,
  !!{function}s કે !!{constant}s આવતાં નથી, એટલે
  $\Struct{M}$નું વ્યાપકક્ષેત્ર જ સંબંધિત છે. હવે
  $s_1(X) = \{1, 2\}$ અને $s_1(Y) = \{3\}$ માટે
  $\Sat{M}{\lforall[z][(\Atom{X}{z}
      \liff \lnot \Atom{Y}{z})]}[s_1]$ મળે છે.

  બીજી તરફ, જો $s_2(X) = \{1, 2\}$ અને $s_2(Y) = \{2, 3\}$ હોય,
  તો $\Sat/{M}{\lforall[z][(\Atom{X}{z} \liff \lnot
      \Atom{Y}{z})]}[s_2]$. કારણ કે
  $\Sat{M}{\Atom{X}{z}}[\Subst{s_2}{2}{z}]$ ($2 \in \Subst{s_2}{2}{z}(X)$ હોવાથી),
  પરંતુ $\Sat/{M}{\lnot \Atom{Y}{z}}[\Subst{s_2}{2}{z}]$
  (કારણ કે $2 \in \Subst{s_2}{2}{z}(Y)$ પણ છે).
\end{ex}

\begin{ex}
  કોઈ $N \subseteq \Domain{M}$ માટે
  $\Sat{M}{(\lexists[y][\Atom{Y}{y}] \land \lforall[z][(\Atom{X}{z}
  \liff \lnot \Atom{Y}{z})])}[\Subst{s}{N}{Y}]$ હોય, તો
  $\Sat{M}{\lexists[Y][(\lexists[y][\Atom{Y}{y}] \land
  \lforall[z][(\Atom{X}{z} \liff \lnot \Atom{Y}{z})])]}[s]$ હોય છે.
  આ માટે $N \neq \emptyset$ હોવું પડે (જેથી
  $\Sat{M}{\lexists[y][\Atom{Y}{y}]}[\Subst{s}{N}{Y}]$ થાય) અને
  અગાઉના ઉદાહરણની જેમ $N = \Domain{M} \setminus s(X)$ હોવું પડે.
  બીજા શબ્દોમાં,
  $\Sat{M}{\lexists[Y][(\lexists[y][\Atom{Y}{y}] \land
  \lforall[z][(\Atom{X}{z} \liff \lnot \Atom{Y}{z})])]}[s]$
  ત્યારે અને માત્ર ત્યારે જ જ્યારે $\Domain{M} \setminus s(X)$
  અરિક્ત હોય, એટલે $s(X) \neq \Domain{M}$. આમ,
  $\Domain{M} = \{1, 2, 3\}$ અને $s(X) = \{1, 2\}$ હોય તો આ
  !!{formula} સંતોષાય છે, પરંતુ $s(X) = \{1, 2, 3\} = \Domain{M}$
  હોય તો સંતોષાતું નથી.

  $s(X) = \Domain{M}$ હોય ત્યારે આ !!{formula} સંતોષાતું નથી, તેથી
  !!{sentence}
  \[
  \lforall[X][\lexists[Y][(\lexists[y][\Atom{Y}{y}] \land
      \lforall[z][(\Atom{X}{z} \liff \lnot \Atom{Y}{z})])]]
  \]
  ક્યારેય સંતોષાતું નથી: કોઈપણ !!{structure}~$\Struct{M}$ માટે
  $s(X) = \Domain{M}$ એવી નિયુક્તિ !!{sentence}ને અસત્ય બનાવશે.
  બીજી બાજુ, વાક્ય
  \[
  \lexists[X][\lexists[Y][(\lexists[y][\Atom{Y}{y}] \land
      \lforall[z][(\Atom{X}{z} \liff \lnot \Atom{Y}{z})])]]
  \]
  કોઈપણ નિયુક્તિ~$s$ને સાપેક્ષ સંતોષાય છે, કારણ કે હંમેશાં
  $M \subseteq \Domain{M}$ અને $M \neq \Domain{M}$ એવો ઉપગણ મળી
  શકે છે (દાખલા તરીકે $M = \emptyset$).
\end{ex}

\begin{ex}
  દ્વિતીય-ક્રમ !!{sentence}~$\lforall[X][\lforall[y][X(y)]]$ કહે છે
  કે દરેક $1$-સ્થાની સંબંધ, એટલે દરેક ગુણધર્મ, દરેક વસ્તુને લાગુ પડે
  છે. તે સ્પષ્ટપણે ક્યારેય સાચું નથી: દરેક~$\Struct{M}$માં
  $s(X) = \emptyset$ અને $s(y) = a \in \Domain{M}$ એવી
  ચલ-નિયુક્તિ~$s$ માટે $\Sat/{M}{X(y)}[s]$ છે. આથી દ્વિતીય-ક્રમ
  તર્કશાસ્ત્રમાં $!A \lif \lforall[X][\lforall[y][X(y)]]$ એ
  $\lnot !A$ને સમકક્ષ છે, એટલે
  $\Sat{M}{!A \lif
  \lforall[X][\lforall[y][X(y)]]}$ ત્યારે અને માત્ર ત્યારે જ
  જ્યારે $\Sat{M}{\lnot !A}$. બીજા શબ્દોમાં, દ્વિતીય-ક્રમ
  તર્કશાસ્ત્રમાં $\lforall$ અને~$\lif$ વડે $\lnot$ વ્યાખ્યાયિત
  કરી શકાય છે.
\end{ex}

\begin{prob}
  બતાવો કે દ્વિતીય-ક્રમ તર્કશાસ્ત્રમાં $\lforall$ અને~$\lif$ વડે
  બીજા સંયોજકો વ્યાખ્યાયિત કરી શકાય છે:
  \begin{enumerate}
    \item સાબિત કરો કે દ્વિતીય-ક્રમ તર્કશાસ્ત્રમાં $!A \land !B$ એ
    $\lforall[X][(!A \lif (!B \lif \lforall[x][X(x)])\lif
    \lforall[x][X(x)])]$ને સમકક્ષ છે.
    \item માત્ર $\lforall$ અને $\lif$ વાપરીને $!A \lor !B$ને
    સમકક્ષ દ્વિતીય-ક્રમ સૂત્ર શોધો.
  \end{enumerate}
\end{prob}

\end{document}
